% ============================================================================
% Methods: how the family amplitudes are estimated (pedagogical version)
% src: family_basis_all.py (span_fit ~L353-357, estimator now always OLS
%      ~L518-522; ridge_amps ~L360-363 kept defined but unused since
%      2026-08-12), lam_fit_amps ~L433-446, sig_pred ~L573), family_model.py;
%      numeric constants recomputed 2026-08-12 from the shipped 8-latent
%      family_model_bundle.npz; robustness numbers from the 2026-08-12
%      bake-off + basis-weighting and anchor-invariance checks (WORKLOG).
% 2026-08-12 (later same day): ridge estimator removed entirely (author
%      decision) -- plain OLS/pinv everywhere, including clk (kappa=520,
%      previously the only statistic to cross the ridge threshold). This
%      fragment was rewritten in place: the "gentle thumb on the scale"
%      section now explains why OLS is safe rather than motivating ridge;
%      the near-twin story and SVD filter analysis are kept as the
%      justification (and as the amplitude-interpretation caveat). See
%      FAMILY_BASIS_METHODS.md's ESTIMATOR addendum for the full numbers.
% Standalone fragment: assumes amsmath; adapt macros to the paper's style.
% ============================================================================

\subsection{How the model is fit: the amplitudes, step by step}
\label{sec:amp-methods}

\paragraph{The model in one sentence.}
Every simulated universe's statistic is written as ``the average universe,
plus a little extra,'' and the little extra is always a mixture of a
handful of fixed, measured shapes:
\begin{equation}
\mathbf{y}(\boldsymbol{\lambda})
\;\simeq\;
\underbrace{\bar{\mathbf{y}}}_{\text{average universe}}
\;+\;
\sum_{k=1}^{K}
\underbrace{a_k(\boldsymbol{\lambda})}_{\text{how much}}\;
\underbrace{\mathbf{b}_k}_{\text{of shape }k}.
\end{equation}
Fitting the model therefore means answering exactly two questions, each of
which turns out to be a small linear regression: \emph{(1) given a
simulated curve, how much of each shape is in it?} and \emph{(2) given
eight numbers measured in a simulation's halos, what mixture should we
expect?}  Everything below is one of those two questions.

\paragraph{The ingredients (all measured, none fit).}
The curves $\mathbf{y}_r$ live on a fixed grid of $P$ bins ($P=24$
$\ell$-bands for $S(\ell)$; $P=22$ $\nu$-bins for the PDF and peak/minima
counts; $P=14$ for the Minkowski functionals).  The \emph{anchor}
$\bar{\mathbf{y}}$ is simply the bin-by-bin average of the $N=256$
Sobol-design curves, and each run's \emph{deviation} is
$\mathbf{d}_r=\mathbf{y}_r-\bar{\mathbf{y}}$.  The \emph{shapes}
$\mathbf{b}_k$ come from the independent twobound experiments: turn one
feedback knob from its low to its high bound, subtract the two curves,
and you get that knob's response shape.  Knobs whose response shapes look
alike (correlation $\ge0.85$) are grouped into a family, and the family's
archetype is the (response-strength-weighted) average of its members'
shapes.  Each archetype is normalized so its largest excursion equals
one.  That normalization is what makes the amplitudes readable:
\emph{$a_k$ is the number of units of shape $k$ present, quoted at shape
$k$'s peak, in the units of the statistic itself} --- for $S(\ell)$,
$a_{C1}=-0.05$ means ``this universe carries a $C_1$-shaped suppression
of depth $0.05$ where $C_1$ peaks.''

\paragraph{Question 1: how much of each shape is in a given curve?}
This is a projection.  We look for the mixture of shapes that lands as
close as possible to the measured deviation:
\begin{equation}
\mathbf{a}_r=\arg\min_{\mathbf{a}}
\Bigl\lVert\, \mathbf{d}_r-\textstyle\sum_k a_k \mathbf{b}_k \,\Bigr\rVert^2 .
\label{eq:amp-objective}
\end{equation}
Least squares has a bookkeeping solution worth seeing once.  Collect the
overlaps between the shapes, $G_{kj}=\mathbf{b}_k\!\cdot\!\mathbf{b}_j$
(the $K\times K$ \emph{Gram matrix} --- its diagonal says how big each
shape is, its off-diagonal how much two shapes resemble each other), and
the overlaps of the data with each shape,
$(\mathsf{B}\mathbf{d}_r)_k=\mathbf{b}_k\!\cdot\!\mathbf{d}_r$.  Then the
best-fit amplitudes solve the small linear system
\begin{equation}
\mathsf{G}\,\mathbf{a}_r=\mathsf{B}\,\mathbf{d}_r .
\label{eq:normal-eq}
\end{equation}
If the shapes were orthogonal, $\mathsf{G}$ would be diagonal and each
amplitude would just be ``overlap of the data with my shape.''  They are
not orthogonal --- and that is where the one subtlety of this section
lives.

\paragraph{The near-twin problem.}
Suppose two shapes were exactly proportional.  Then only the \emph{sum}
of their amplitudes would matter: $(a_1,a_2)=(10.3,-10)$ and $(0.2,0.1)$
would produce identical curves, and the fit could return either,
depending on noise.  Our families are not identical twins, but for
$S(\ell)$ they are close: the four archetypes are smooth suppression-like
curves that substantially resemble one another.  The degree of
near-twinness is measured by the conditioning number
$\kappa=\sigma_1/\sigma_K$ (the ratio of the largest to smallest singular
value of $\mathsf{B}$): $\kappa=1$ means perfectly independent shapes,
large $\kappa$ means some \emph{combination} of shapes is nearly
invisible.  Measured values:
\begin{center}
\begin{tabular}{lccl}
statistic & $\kappa(\mathsf{B})$ & estimator & singular values of $\mathsf{B}$\\
\hline
$S(\ell)$    & 520  & OLS & $(4.57,\,2.18,\,0.216,\,0.0088)$\\
PDF          & 91   & OLS   & $(2.46,\,1.04,\,0.68,\,0.027)$\\
$N_{\rm pk}$ & 10   & OLS   & $(3.41,\,1.35,\,1.02,\,0.61,\,0.33)$\\
$N_{\rm min}$& 7.3  & OLS   & $(2.68,\,0.89,\,0.60,\,0.37)$\\
$V_0$        & 49   & OLS   & $(2.89,\,0.63,\,0.060)$\\
$V_1$        & 5.7  & OLS   & $(2.52,\,0.44)$\\
$V_2$        & 1    & OLS   & $(1.93)$\\
\end{tabular}
\end{center}
Every statistic --- $S(\ell)$ included --- uses the same plain-least-squares
estimator (Eq.~\ref{eq:normal-eq} with $\alpha=0$; the minimum-norm
solution when $\mathsf{G}$ is singular). $S(\ell)$'s $\kappa=520$ used to
trigger a special ridge branch here; the next paragraph explains why that
branch was removed and why doing so is safe.

\paragraph{Why plain OLS is safe here.}
A tempting fix for a near-twin basis is to add a small penalty on the
amplitudes themselves, $+\,\alpha\lVert\mathbf{a}\rVert^2$ in
Eq.~\ref{eq:amp-objective} (``ridge''), turning
Eq.~\ref{eq:normal-eq} into
$(\mathsf{G}+\alpha\mathbb{1})\,\mathbf{a}_r=\mathsf{B}\,\mathbf{d}_r$ ---
in words, \emph{among mixtures that fit the curve almost equally well,
prefer the one with small weights}. $S(\ell)$ shipped with exactly this
fix for a time ($\alpha=10^{-2}\,\mathrm{tr}(\mathsf{G})/K=0.064$). We
have since removed it in favor of one estimator everywhere, and the SVD
filter-factor analysis that motivated ridge is precisely what shows the
removal costs nothing.

The singular value decomposition sorts amplitude space into directions
ranked by how visible they are in the curve: direction $j$ produces a
curve of size $\sigma_j$ per unit amplitude. Ridge multiplies the fitted
coefficient of each direction by a filter
$f_j=\sigma_j^2/(\sigma_j^2+\alpha)$; for $S(\ell)$ at the $\alpha$ above,
\[
f=(0.997,\;0.987,\;0.42,\;0.0012),
\]
so ridge would leave directions 1--2 (visible) essentially untouched,
half-trust direction 3, and freeze direction 4. Direction 4 deserves its
number: it is the amplitude combination
$\mathbf{u}_4=(-0.76,+0.59,+0.27,+0.09)$ of $(C_1,\dots,C_4)$, and one
full unit of it changes the predicted curve by at most $0.005$ --- half a
percent of a unit-peak shape, below the model's own predictive
uncertainty everywhere on the grid. That is exactly why plain least
squares is safe: setting $\alpha=0$ sets every filter factor to $1$,
including $f_4$, so it un-freezes $\mathbf{u}_4$ and assigns it a
coefficient amplified by $1/\sigma_4\simeq115$ --- but because $\mathbf{u}_4$'s
curve image is observationally invisible, that amplification changes the
\emph{amplitude values} without moving the \emph{reconstructed curve}
(span $R^2$ goes from $0.9984$ under ridge to $0.9997$ under OLS ---
OLS reaches the free-amplitude ceiling exactly, ridge was $1.3\times10^{-3}$
\emph{below} it) or the model's end-to-end accuracy (measured below). The
estimator choice is therefore purely an \emph{identifiability} matter, not
an accuracy trade-off, and the safest way to remove a special case is to
remove it. Two habits follow, unchanged from when ridge existed: compare
amplitudes between data sets only when both were fit with the \emph{same}
estimator (mixing conventions --- e.g.\ one plain-least-squares point
plotted against a ridge-fitted cloud --- lands hundreds of ``sigma'' away
along $\mathbf{u}_4$ while representing a nearly identical curve; this
repository shipped that exact bug once, in the amplitude corner plot,
before both were unified onto OLS), and interpret \emph{individual}
amplitudes only through the well-measured directions of $\mathsf{B}$
--- never $\mathbf{u}_4$ alone.

\paragraph{Two facts we get for free.}
(1) The amplitude fit is a fixed matrix applied to each run's own curve
($\mathsf{A}=\mathsf{D}\mathsf{X}$, one row at a time), so run $r$'s
amplitudes never see any other run.  Cross-validation of the next step is
therefore clean by construction.  (2) The amplitudes average to exactly
zero over the design --- deviations from an average always do
($\overline{\mathsf{A}}=\overline{\mathsf{D}}\mathsf{X}=\mathbf{0}$;
verified at the $10^{-16}$ level).  This is why the model needs no fitted
constant term below.

\paragraph{Question 2: predicting the mixture from eight halo numbers.}
The latents $\boldsymbol{\lambda}$ are eight medians measured in a
simulation's halo population (baryon fraction, stellar fraction, gas
concentration, temperature in a group-mass bin; electron pressure and
$R_{200c}$ gas fraction in the group and cluster bins ---
Table~\ref{tab:latents}).  They are measurements, never fit parameters.
For each amplitude we now fit \emph{one straight line in eight
variables}:
\begin{equation}
a_k(\boldsymbol{\lambda})
= \mathbf{M}_k\cdot(\boldsymbol{\lambda}-\boldsymbol{\lambda}_{\rm ref}),
\qquad
\boldsymbol{\lambda}_{\rm ref}=\tfrac1N\textstyle\sum_r\boldsymbol{\lambda}_r ,
\label{eq:latmap}
\end{equation}
with the slopes $\mathsf{M}$ obtained by ordinary least squares over the
256 design runs.  No intercept appears because none is needed: the
latents are centered on their design mean and the amplitudes average to
zero, so the line passes through the origin of both.  (The bundle stores
the intercept anyway; it is identically zero.)

One interpretive warning, because the design is not a grid: the eight
latents move \emph{together} across the design --- universes that expel
gas also puff it up ($r(\tilde f_{\rm bar},c_{\rm gas})=0.94$).  The rows
of $\mathsf{M}$ are \emph{partial} slopes --- the response to changing
one latent while holding the other seven fixed --- and these are not the
slopes you would read off a scatter plot of $a_k$ against one latent
alone (those ``marginal'' slopes fold in every co-moving latent).  The
construction figure draws exactly this distinction, and the tutorial's
partial-vs-conditional sweeps are its curve-level counterpart.

\paragraph{Putting it together.}
Chaining the two regressions, the whole model is one affine map from
measured halo numbers to curves:
\begin{equation}
\hat{\mathbf{y}}(\boldsymbol{\lambda})
= \bar{\mathbf{y}}
+ \Theta\,(\boldsymbol{\lambda}-\boldsymbol{\lambda}_{\rm ref}),
\qquad
\Theta=\mathsf{B}^{\!\top}\mathsf{M}^{\!\top}\in\mathbb{R}^{P\times L}.
\label{eq:composed}
\end{equation}
$\Theta$ is the model's full answer to ``how does bin $p$ respond to
latent $j$'' --- factored into a shape dictionary measured from one-knob
experiments ($\mathsf{B}$) and a sensitivity matrix fit on the design
($\mathsf{M}$).  The anchor cannot be absorbed into the shapes: the
$\mathbf{b}_k$ are built from response \emph{differences}, so they span
the directions in which curves \emph{change}, while $\bar{\mathbf{y}}$
records where the population \emph{sits} ($8.9\%$ of $\bar S-1$ lies
outside the shapes' span --- $5\times10^{-3}$ in $S$ units, larger than
the model's own precision at low $\ell$).

\paragraph{How good is it, honestly.}
Hide one fifth of the runs; refit the line (only the line --- the
amplitudes are per-run, see above); predict the hidden curves from their
measured latents; repeat five times so every run is hidden once.  The
pooled score
$R^2_{\rm e2e}=1-\sum_r\lVert\mathbf{d}_r-\hat{\mathbf{d}}_r\rVert^2
/\sum_r\lVert\mathbf{d}_r\rVert^2$ is what we quote per statistic, and
the per-bin scatter of those held-out errors,
$\sigma_{\rm pred}(x_p)=\mathrm{std}_r[d_{rp}-\hat d^{\rm CV}_{rp}]$, is
the shaded band drawn around every prediction (it matches a full
leave-one-out computation to $\sim$1\%).  Read the band for what it is:
an \emph{in-design average}.  It is strongly scale-dependent for
$S(\ell)$ --- $0.004$ at $\ell\sim1.5\times10^{3}$, where the latents pin
the response, versus $0.019$ at the $\ell\sim1.3\times10^{4}$ trough,
where eight halo numbers genuinely underdetermine the small-scale
physics.

\paragraph{Stress tests (measured, not asserted).}
Choices that sound consequential and are not: replacing the
response-strength weighting in the family archetypes by a plain mean
moves the shapes by $<0.5\%$ and $R^2_{\rm e2e}$ by at most $0.011$
(PDF, whose 16-member family has many weak, noisy members --- the
weighting is an approximate inverse-variance weight); the $S(\ell)$
member responses being ratios to a fiducial reference rather than to the
DMO leaves a smooth $\le13.5\%$ reweighting on the shapes that the
amplitudes absorb entirely (span changes by $3\times10^{-5}$);
re-anchoring the model on any curve within the shapes' span of
$\bar{\mathbf{y}}$ --- the measured fiducial curve included --- is a pure
relabeling of amplitudes that moves predictions by
$\le1.3\times10^{-3}$; and, most directly, \emph{removing the ridge
estimator itself} for $S(\ell)$ (the one statistic that ever used it) ---
measured, not asserted, on the pipeline's own end-to-end cross-validation:
$R^2_{\rm e2e}=0.9585$ (ridge) vs.\ $0.9593$ (OLS), and the fiducial
trough prediction (the single hardest number this model produces, band
$\ell\approx1.3\times10^4$) moves from $0.9112$ to $0.908$ against the
same $0.881$ measured value and the same $\sim0.019$ predictive $\sigma$
--- both within $1.5\sigma$ of measurement, neither reading as a
regression. This is the measured version of the filter-factor argument
above: the estimator moves the amplitudes along $\mathbf{u}_4$ and leaves
everything the model is actually graded on --- the curve --- alone.
